\Needspace{14\baselineskip}
\Article{30}

Eine weitere Methode, die Summe $W$ allgemein zu bestimmen, ergibt sich aus unseren Ausführungen in den Art. $22.\ 24$.\quad Setzen wir $\cos\omega+i\sin\omega=\rho$ und

\[
\begin{aligned}
\rho^{\frac{nn}{aa}}=\alpha,\quad\rho^{\frac{nn}{bb}}=\beta,\quad\rho^{\frac{nn}{cc}}=\gamma\;\text{usw.}
\end{aligned}
\]

so dass $r=\rho^k$, $A=\alpha^k$, $B=\beta^k$, $C=\gamma^k$ usw. gilt.\quad Dann ist

\[
\begin{aligned}
1+\rho+\rho^4+\rho^9+\text{usw.}+\rho^{(n-1)^2}
\end{aligned}
\]

das Produkt der Faktoren

\[
\begin{aligned}
&1+\alpha+\alpha^4+\alpha^9+\text{usw.}+\alpha^{(a-1)^2}\\
&1+\beta+\beta^4+\beta^9+\text{usw.}+\beta^{(b-1)^2}\\
&1+\gamma+\gamma^4+\gamma^9+\text{usw.}+\gamma^{(c-1)^2},\;\text{usw.}
\end{aligned}
\]

\SourceLeaf{50}

\Needspace{12\baselineskip}
und somit $W$ das Produkt der Faktoren

\[
\begin{aligned}
w&=1+\rho+\rho^4+\rho^9+\text{usw.}+\rho^{(n-1)^2}\\[2pt]
\mathfrak A&=\frac{1+A+A^4+A^9+\text{usw.}+A^{(a-1)^2}}{1+\alpha+\alpha^4+\alpha^9+\text{usw.}+\alpha^{(a-1)^2}}\\[2pt]
\mathfrak B&=\frac{1+B+B^4+B^9+\text{usw.}+B^{(b-1)^2}}{1+\beta+\beta^4+\beta^9+\text{usw.}+\beta^{(b-1)^2}}\\[2pt]
\mathfrak C&=\frac{1+C+C^4+C^9+\text{usw.}+C^{(c-1)^2}}{1+\gamma+\gamma^4+\gamma^9+\text{usw.}+\gamma^{(c-1)^2}}
\end{aligned}
\]

Der erste Faktor $w$ ist bereits durch die oben angegebenen Untersuchungen bestimmt (Art. $19$); die übrigen Faktoren $\mathfrak A,\mathfrak B,\mathfrak C$ usw. ergeben sich aus den Formeln der Art. $22.\ 24$, die wir hier nochmals zusammenstellen, um alles beisammen zu haben\begingroup\renewcommand{\thefootnote}{*)}\footnote{Was dort $k$ und $h$ war, ist hier offenbar $\frac{n}{a}$ und $k$ für den zweiten Faktor, $\frac{n}{b}$ und $k$ für den dritten Faktor usw.}\endgroup.\quad Hier sind zwölf Fälle zu unterscheiden, nämlich

I.\quad Ist $a$ eine (ungerade) Primzahl $=p$ oder eine Potenz einer solchen Zahl mit ungeradem Exponenten und $k$ ein quadratischer Rest von $p$, so ist der entsprechende Faktor $\mathfrak A=+1$.

II.\quad Bleiben die übrigen Bedingungen unverändert und ist $k$ ein quadratischer Nichtrest von $p$, so ist $\mathfrak A=-1$.

III.\quad Ist $a$ das Quadrat einer ungeraden Primzahl oder eine höhere Potenz derselben mit geradem Exponenten, so ist $\mathfrak A=+1$.

IV.\quad Ist $a$ $=4$ oder eine höhere Zweierpotenz mit geradem Exponenten und hat zugleich $k$ die Form $4\mu+1$, so ist $\mathfrak A=+1$.

V.\quad Sind die übrigen Bedingungen wie in IV, hat $k$ die Form $4\mu+3$ und $\frac{n}{a}$ die Form $4\mu+1$, so ist $\mathfrak A=-i$.

VI.\quad Sind die übrigen Bedingungen wie in IV, hat $k$ die Form $4\mu+3$ und $\frac{n}{a}$ die Form $4\mu+3$, so ist $\mathfrak A=+i$.

VII.\quad Ist $a$ $=8$ oder eine höhere Zweierpotenz mit ungeradem Exponenten und hat $k$ die Form $8\mu+1$, so ist $\mathfrak A=+1$.

VIII.\quad Sind die übrigen Bedingungen wie in VII und hat $k$ die Form $8\mu+5$, so ist $\mathfrak A=-1$.

IX.\quad Sind die übrigen Bedingungen wie in VII, hat $k$ die Form $8\mu+3$ und $\frac{n}{a}$ die Form $4\mu+1$, so ist $\mathfrak A=+i$.

\SourceLeaf{51}

X.\quad Sind die übrigen Bedingungen wie in VII, hat $k$ die Form $8\mu+3$ und $\frac{n}{a}$ die Form $4\mu+3$, so ist $\mathfrak A=-i$.

XI.\quad Sind die übrigen Bedingungen wie in VII, hat $k$ die Form $8\mu+7$ und $\frac{n}{a}$ die Form $4\mu+1$, so ist $\mathfrak A=-i$.

XII.\quad Sind die übrigen Bedingungen wie in VII, hat $k$ die Form $8\mu+7$ und $\frac{n}{a}$ die Form $4\mu+3$, so ist $\mathfrak A=+i$.

Den Fall $a=2$ übergehen wir; hier wäre $\mathfrak A$ zwar $=\frac00$, also unbestimmt, doch gilt dann stets $W=0$.

Die übrigen Faktoren $\mathfrak B,\mathfrak C$ usw. hängen von $b,c$ usw. ebenso ab wie $\mathfrak A$ von $a$, soweit diese Zahlen in ihre Bestimmung eingehen.

\Needspace{14\baselineskip}
\Article{31}

Nach dieser zweiten Methode gestaltet sich das erste Beispiel aus Art. $29$ folgendermaßen:

{\small
\[
\begin{aligned}
&\text{Der Faktor }w\text{ ist }=(1+i)\surd2520\\
&\text{Für }a=8\text{ ist der zugehörige Faktor }\mathfrak A\text{ nach Fall VIII }=-1\\
&\text{Zum zweiten Faktor von }n\text{, }9\text{, gehört der Faktor }+1\;\text{(nach Fall III.)}\\
&\text{Zum Faktor }5\text{ gehört der Faktor }-1\;\text{(nach Fall II.)}\\
&\text{Zum Faktor }7\text{ gehört der Faktor }-1\;\text{(nach Fall II.)}
\end{aligned}
\]
}

Daraus ergibt sich das Produkt $W=(-1-i)\surd2520$, wie in Art. $29$.

\Needspace{18\baselineskip}
\Article{32}

Da sich der Wert von $W$ nach \textit{zwei} Methoden bestimmen lässt, von denen die eine auf den Beziehungen der Zahlen $\frac{nk}{a}$, $\frac{nk}{b}$, $\frac{nk}{c}$ usw. zu den Zahlen $a,b,c$ usw. beruht, die andere dagegen von den Beziehungen von $k$ zu den Zahlen $a,b,c$ usw. abhängt, muss zwischen all diesen Beziehungen ein bedingter Zusammenhang bestehen, so dass jede von ihnen aus den übrigen bestimmbar sein muss.\quad Nehmen wir an, alle Zahlen $a,b,c$ usw. seien ungerade Primzahlen, und setzen wir $k$ $=1$. Wir teilen die Faktoren $a,b,c$ usw. in zwei Klassen: Die eine enthält diejenigen von der Form $4\mu+1$, die mit $p,p',p''$ usw. bezeichnet werden; die andere besteht aus denen von der Form $4\mu+3$, die wir mit $q,q',q''$ usw. bezeichnen. Die Anzahl der Letzteren bezeichnen wir mit $m$. Dabei bemerken wir zunächst, dass $n$ die Form $4\mu+1$ hat, wenn $m$ gerade ist (hierzu gehört auch der Fall, dass überhaupt keine Faktoren der zweiten Klasse vorhanden sind, also $m=0$), dagegen $n$ die Form $4\mu+3$ hat, wenn $m$ ungerade ist.\quad Die Bestimmung\SourceLeaf{52} von $W$ nach der ersten Methode wird nun folgendermaßen ausgeführt.\quad Die Zahlen $P,P'\;P''$ usw., $Q,Q',Q''$ usw. sollen von den Beziehungen der Zahlen $\frac{n}{p}$, $\frac{n}{p'}$, $\frac{n}{p''}$ usw., $\frac{n}{q}$, $\frac{n}{q'}$, $\frac{n}{q''}$ usw. zu den Zahlen $p,p',p''$ usw., $q,q',q''$ usw. jeweils so abhängen, dass wir setzen:

\[
\begin{aligned}
P&=+1,\;\text{wenn }\frac{n}{p}\text{ ein quadratischer Rest von }p\text{ ist}\\
P&=-1,\;\text{wenn }\frac{n}{p}\text{ ein quadratischer Nichtrest von }p\text{ ist}
\end{aligned}
\]

und ebenso für die übrigen.\quad Dann ist $W$ das Produkt der Faktoren $P\surd p$, $P'\surd p'$, $P''\surd p''$ usw. $iQ\surd q$, $iQ'\surd q'$, $iQ''\surd q''$ usw., also

\[
\begin{aligned}
W=PP'P''\ldots QQ'Q''\ldots i^m\surd n
\end{aligned}
\]

Nach der zweiten Methode oder vielmehr unmittelbar nach den Regeln des Art. $19$ gilt

\[
\begin{aligned}
W&=+\surd n,\;\text{wenn }n\text{ die Form }4\mu+1\text{ hat},\\&\qquad\text{oder gleichbedeutend, wenn }m\text{ gerade ist}\\
W&=+i\surd n,\;\text{wenn }n\text{ die Form }4\mu+3\text{ hat},\\&\qquad\text{oder wenn }m\text{ ungerade ist.}
\end{aligned}
\]

Beide Fälle lassen sich gemeinsam durch die folgende Formel ausdrücken:

\[
\begin{aligned}
W=i^{mm}\surd n
\end{aligned}
\]

Daraus folgt also

\[
\begin{aligned}
PP'P''\ldots QQ'Q''\ldots=i^{mm-m}
\end{aligned}
\]

Nun ist $i^{mm-m}$ $=1$, wenn $m$ die Form $4\mu$ oder $4\mu+1$ hat, und $=-1$, wenn $m$ die Form $4\mu+2$ oder $4\mu+3$ hat. Daraus leiten wir den folgenden höchst eleganten Satz ab:

\textsc{Satz.}\quad \textit{Es bezeichnen $a,b,c$ usw. verschiedene positive ungerade Primzahlen, deren Produkt $=n$ gesetzt wird; unter ihnen seien $m$ von der Form $4\mu+3$, die übrigen von der Form $4\mu+1$. Die Anzahl derjenigen unter diesen Zahlen $a,b,c$ usw., von denen jeweils $\frac{n}{a}$, $\frac{n}{b}$, $\frac{n}{c}$ usw. Nichtreste sind, ist gerade, wenn $m$ die Form $4\mu$ oder $4\mu+1$ hat, dagegen ungerade, wenn $m$ die Form $4\mu+2$ oder $4\mu+3$ hat.}

Setzen wir beispielsweise $a=3$, $b=5$, $c=7$, $d=11$, so haben wir drei Zahlen von der Form $4\mu+3$, nämlich $3$, $7$ und $11$; dabei gilt $5\pdot7\pdot11\,R3$; $3\pdot7\pdot11\,R5$; $3\pdot5\pdot11\,R7$; $3\pdot5\pdot7\,N11$, also ist allein $\frac{n}{d}$ ein Nichtrest von $d$.

\Needspace{12\baselineskip}
\Article{33}

Das berühmte \textit{Fundamentaltheorem} über quadratische Reste ist nichts anderes als ein besonderer Fall des soeben entwickelten Satzes.\quad Beschränken wir nämlich die Anzahl\SourceLeaf{53} der Zahlen $a,b,c$ usw. auf \textit{zwei}, so müssen, wenn nur eine von ihnen oder keine die Form $4\mu+3$ hat, offenbar entweder zugleich $aRb$, $bRa$ oder zugleich $aNb$, $bNa$ gelten; haben dagegen beide die Form $4\mu+3$, so muss die eine ein Nichtrest der anderen und diese ein Rest jener sein.\quad Hier liegt also ein \textit{vierter} Beweis dieses höchst wichtigen Satzes vor, dessen ersten und zweiten Beweis wir in den Disquisitiones Arithmeticae und dessen dritten wir kürzlich in einer besonderen Abhandlung gegeben haben (\textit{Commentt. T. XVI}); zwei weitere, die wiederum auf völlig anderen Grundsätzen beruhen, werden wir später darlegen.\quad Es ist wahrlich höchst erstaunlich, dass dieser so schöne Satz, der anfangs allen Versuchen so hartnäckig entgangen war, später auf so vielen, himmelweit voneinander entfernten Wegen erreicht werden konnte.

\Needspace{14\baselineskip}
\Article{34}

Auch die übrigen Sätze, die gewissermaßen eine Ergänzung zum Fundamentaltheorem bilden, nämlich diejenigen, durch die man die Primzahlen erkennt, von denen $-1,+2$ und $-2$ Reste oder Nichtreste sind, lassen sich aus denselben Grundsätzen ableiten. Wir beginnen mit dem Rest $+2$.

Setzen wir $n=8a$, wobei $a$ eine Primzahl ist, und $k=1$, so ist nach der Methode des Art. $28$ $W$ das Produkt zweier Faktoren. Der eine ist $+\surd a$ oder $+i\surd a$, wenn $8$ oder gleichbedeutend $2$ ein quadratischer Rest von $a$ ist; dagegen $-\surd a$ oder $-i\surd a$, wenn $2$ ein Nichtrest von $a$ ist.\quad Der zweite Faktor ist

\[
\begin{aligned}
&(1+i)\surd8,\quad\text{wenn }a\text{ die Form }8\mu+1\text{ hat}\\
&(-1+i)\surd8,\quad\text{wenn }a\text{ die Form }8\mu+3\text{ hat}\\
&(-1-i)\surd8,\quad\text{wenn }a\text{ die Form }8\mu+5\text{ hat}\\
&(1-i)\surd8,\quad\text{wenn }a\text{ die Form }8\mu+7\text{ hat}
\end{aligned}
\]

Nach Art. $18$ ist aber stets $W=(1+i)\surd n$; teilt man diesen Wert durch die vier Werte des zweiten Faktors, so erkennt man, dass der erste Faktor

\[
\begin{aligned}
&+\surd a,\quad\text{wenn }a\text{ die Form }8\mu+1\text{ hat}\\
&-i\surd a,\quad\text{wenn }a\text{ die Form }8\mu+3\text{ hat}\\
&-\surd a,\quad\text{wenn }a\text{ die Form }8\mu+5\text{ hat}\\
&+i\surd a,\quad\text{wenn }a\text{ die Form }8\mu+7\text{ hat}
\end{aligned}
\]

sein muss. Daraus folgt unmittelbar, dass $2$ im ersten und vierten Fall ein Rest von $a$, im zweiten und dritten Fall dagegen ein Nichtrest sein muss.

\SourceLeaf{54}

\Needspace{14\baselineskip}
\Article{35}

Die Primzahlen, von denen $-1$ ein Rest oder Nichtrest ist, lassen sich mithilfe des folgenden Satzes leicht erkennen, der auch für sich genommen recht bemerkenswert ist.

\textsc{Satz.}\quad \textit{Das Produkt der beiden Faktoren}

\[
\begin{aligned}
W'&=1+r^{-1}+r^{-4}+\text{usw.}+r^{-(n-1)^2}\\
W&=1+r+r^4+\text{usw.}+r^{(n-1)^2}
\end{aligned}
\]

\textit{ist $=n$, wenn $n$ ungerade ist; oder $=0$, wenn $n$ das Doppelte einer ungeraden Zahl ist; oder $=2n$, wenn $n$ durch vier teilbar ist.}

\Needspace{12\baselineskip}
\textit{Beweis.}\quad Da offenbar

\[
\begin{aligned}
W&=r+r^4+r^9+\text{usw.}+r^{nn}\\
&=r^4+r^9+\text{usw.}+r^{(n+1)^2}\\
&=r^9+\text{usw.}+r^{(n+2)^2}\quad\text{usw.}
\end{aligned}
\]

gilt, lässt sich das Produkt $WW'$ auch folgendermaßen darstellen:

\[
\begin{aligned}
&1+r+r^4+r^9+\text{usw.}+r^{(n-1)^2}\\
+{}&r^{-1}(r+r^4+r^9+r^{16}+\text{usw.}+r^{nn})\\
+{}&r^{-4}(r^4+r^9+r^{16}+r^{25}+\text{usw.}+r^{(n+1)^2})\\
+{}&r^{-9}(r^9+r^{16}+r^{25}+r^{36}+\text{usw.}+r^{(n+2)^2})\\
&\text{usw.}\\
+{}&r^{-(n-1)^2}(r^{(n-1)^2}+r^{nn}+r^{(n+1)^2}\\&\qquad+r^{(n+2)^2}+\text{usw.}+r^{(2n-2)^2})
\end{aligned}
\]

Summiert man diesen Ausdruck senkrecht, so erhält man

\[
\begin{aligned}
&n\\
+{}&r(1+rr+r^4+r^6+\text{usw.}+r^{2n-2})\\
+{}&r^4(1+r^4+r^8+r^{12}+\text{usw.}+r^{4n-4})\\
+{}&r^9(1+r^6+r^{12}+r^{18}+\text{usw.}+r^{6n-6})\\
+{}&\text{usw.}\\
+{}&r^{(n-1)^2}(1+r^{2n-2}+r^{4n-4}\\&\qquad+r^{6n-6}+\text{usw.}+r^{2(n-1)^2})
\end{aligned}
\]

Ist nun $n$ ungerade, so sind die einzelnen Teile dieses Ausdrucks außer dem ersten, $n$, sämtlich $=0$; denn der zweite ist offenbar $\frac{r(1-r^{2n})}{1-rr}$, der dritte $\frac{r^4(1-r^{4n})}{1-r^4}$ usw. Ist dagegen $n$ gerade, so muss außerdem der Teil

\SourceLeaf{55}

\[
\begin{aligned}
r^{\frac14nn}(1+r^n+r^{2n}+r^{3n}+\text{usw.}+r^{nn-n})
\end{aligned}
\]

ausgenommen werden, der $=nr^{\frac14nn}$ ist. Im ersten Fall ist daher $WW'=n$, im zweiten dagegen $=n+nr^{\frac14nn}$; nun ist $r^{\frac14nn}$ $=+1$, wenn $n$ durch vier teilbar ist, und dann ergibt sich $WW'=2n$; dagegen gilt $r^{\frac14nn}=-1$, wenn $n$ das Doppelte einer ungeraden Zahl ist, so dass dann $WW'=0$ herauskommt. Q. E. D.

\Needspace{10\baselineskip}
\Article{36}

Nach Art. 22 steht fest, dass für eine ungerade Primzahl $n$ der Quotient $\frac{W'}{W}$ $=+1$ oder $=-1$ ist, je nachdem, ob $-1$ ein Rest oder Nichtrest von $n$ ist. Daher muss im ersten Fall $W^2=+n$, im zweiten $W^2=-n$ gelten; nach Art. 13 schließen wir somit, dass der erste Fall nur eintreten kann, wenn $n$ die Form $4\mu+1$ hat, und der zweite, wenn $n$ die Form $4\mu+3$ hat.

Aus der Verbindung der für die Reste $+2$ und $-1$ gefundenen Bedingungen folgt schließlich unmittelbar, dass $-2$ ein Rest jeder Primzahl von der Form $8\mu+1$ oder $8\mu+3$ und ein Nichtrest jeder Primzahl von der Form $8\mu+5$ oder $8\mu+7$ ist.

\BookRule{42mm}
